\documentclass[conference]{corl_2026}
\usepackage{amsmath,amssymb,amsfonts}
\usepackage{booktabs}
\usepackage{multirow}
\usepackage{graphicx}
\usepackage{xcolor}
\usepackage{hyperref}
\usepackage{algorithm}
\usepackage{algorithmic}

\usepackage{amsthm}
\newtheorem{proposition}{Proposition}

\newcommand{\ours}{IMM}
\newcommand{\acis}{ACIS}
\newcommand{\deltak}{\Delta_k}

\title{Importance Margin Maximization: A Robustness-Aware Framework\\for Vision-Language-Action Token Compression}

\author{
  Anonymous Authors
}

\begin{document}
\maketitle

%==============================================================================
\begin{abstract}
%==============================================================================
Token compression is essential for deploying Vision-Language-Action (VLA) models on resource-constrained robots, yet existing compression methods are brittle: small input perturbations destabilize importance rankings and cause safety-critical tokens to be discarded.
We introduce \emph{Importance Margin Maximization} (\ours{}), the first robustness-aware compression framework for VLAs with theoretical grounding.
\ours{} trains a lightweight Adaptive Compression Importance Scorer (\acis{}) to maximize the importance margin $\deltak$ between retained and pruned tokens, providing a principled design criterion: $P(\text{top-}k\text{ flip}) \leq k \cdot \varepsilon \cdot L / \deltak$.
A two-stage training procedure decouples \acis{} from the backbone to prevent co-adaptation, achieving near-lossless 50\% token compression with strong robustness guarantees.
On LIBERO-Plus with 7-dimension visual perturbations, \ours{}-compressed SmolVLA recovers 78\% of the robustness lost by standard compression (37.3\% $\rightarrow$ 50.3\% vs.\ 53.1\% uncompressed), while achieving 1.62$\times$ inference speedup.
\ours{} is compression-agnostic: applied as a plugin to FastV, VLA-IAP, and VLA-Cache, it improves perturbation preservation by +53pp on average with negligible overhead.
\end{abstract}

%==============================================================================
\section{Introduction}
\label{sec:intro}
%==============================================================================

Vision-Language-Action (VLA) models enable general-purpose robot manipulation by leveraging pretrained vision-language backbones for action prediction. However, the computational cost of processing hundreds of visual tokens at each control step poses a fundamental bottleneck for real-time deployment.
Token compression---retaining only the most informative visual tokens---has emerged as the dominant approach, with methods such as FastV~\cite{fastv}, VLA-IAP~\cite{vlaiap}, and VLA-Cache~\cite{vlacache} achieving 50--70\% compression ratios with minimal clean-condition degradation.

However, recent work reveals a critical vulnerability: token compression amplifies adversarial fragility. ``Less Is More Until It Breaks''~\cite{lessismore} demonstrates that small input perturbations---camera jitter, lighting changes, or adversarial noise---destabilize importance rankings, causing compressed VLAs to discard safety-critical tokens. The root cause is that existing importance scorers produce rankings with negligible margins between retained and pruned tokens: a perturbation of magnitude $\varepsilon$ can flip the top-$k$ selection when the importance gap $\deltak$ is small.

We identify that compression robustness depends on the \emph{importance margin} $\deltak$, not on token entropy or attention magnitude. This insight leads to a principled design criterion: directly maximizing $\deltak$ provably reduces the probability of selection instability. Concretely, for a scorer with Lipschitz constant $L$, we derive:
\begin{equation}
P(\text{top-}k \text{ set changes under } \|\delta\| \leq \varepsilon) \leq \frac{k \cdot \varepsilon \cdot L}{\deltak}
\label{eq:bound}
\end{equation}
This bound motivates our training objective and provides a principled criterion for evaluating compression robustness.

Our framework, \emph{Importance Margin Maximization} (\ours{}), consists of three components: (1) an Adaptive Compression Importance Scorer (\acis{}) trained to maximize $\deltak$; (2) a two-stage training procedure that decouples \acis{} from the VLA backbone, preventing co-adaptation that would damage backbone quality; and (3) a deployment pipeline where \acis{} scores tokens, applies top-$k$ selection, and feeds the compressed sequence to the frozen backbone.

A key engineering finding is that \acis{} and the backbone \emph{must} be trained in separate stages. Joint training causes the backbone to co-adapt to \acis{}'s mask pattern, creating a degenerate dependency where the backbone cannot function without \acis{} (Section~\ref{sec:ablation_training}). Our two-stage design eliminates this pathology entirely.

We validate \ours{} on SmolVLA (0.5B) and OpenVLA (7B) across multiple benchmarks (aloha\_sim\_insertion, LIBERO-Object/Spatial/Goal, LIBERO-Plus). Our contributions:
\begin{itemize}
\item The first robustness-aware token compression framework for VLAs, grounded in a formal stability condition that serves as a principled design criterion for training objectives.
\item A two-stage training procedure achieving near-lossless 50\% compression with robustness guarantees, Pareto-dominating standard compression on the robustness-efficiency frontier.
\item Demonstration that \ours{} is compression-agnostic, improving robustness by +53pp when applied as a plugin to three existing compression methods.
\end{itemize}

%==============================================================================
\section{Related Work}
\label{sec:related}
%==============================================================================

\paragraph{VLA Token Compression.}
A growing body of work addresses the computational cost of VLA deployment through token reduction. FastV~\cite{fastv} uses attention-weighted token selection; VLA-IAP~\cite{vlaiap} learns input-adaptive pruning policies; VLA-Cache~\cite{vlacache} applies KV-cache eviction. These methods focus exclusively on efficiency under clean conditions and provide no robustness analysis.

\paragraph{Adversarial Robustness for VLAs.}
RobustVLA~\cite{robustvla} applies Jacobian regularization to reduce sensitivity but does not address compression. ProGAL-VLA proposes perturbation-aware fine-tuning without considering token selection stability. Neither line of work connects robustness to the compression mechanism itself.

\paragraph{Certified Robustness.}
Randomized smoothing~\cite{cohen2019certified} provides certified $\ell_2$ robustness but requires inference-time noise injection that degrades clean performance. Our bound (Eq.~\ref{eq:bound}) instead certifies the \emph{selection mechanism}, requiring no inference-time modification.

%==============================================================================
\section{Method}
\label{sec:method}
%==============================================================================

\subsection{Problem Formulation}

Consider a VLA model $f_\theta$ processing $n$ visual tokens $\{v_1, \ldots, v_n\}$. A compression function $C_k: \mathbb{R}^{n \times d} \rightarrow \mathbb{R}^{k \times d}$ selects the top-$k$ tokens by importance score $s_i = g_\phi(v_i, \mathbf{v})$, where $g_\phi$ is a learned scorer. We define the \emph{importance margin}:
\begin{equation}
\deltak = s_{(k)} - s_{(k+1)}
\end{equation}
where $s_{(i)}$ denotes the $i$-th largest score. When $\deltak$ is small, an input perturbation $\|\delta\| \leq \varepsilon$ can change the sorted order at the $k$-th position, altering which tokens are retained.

\subsection{Formal Stability Condition}
\label{sec:formal}

\begin{proposition}
Let $g_\phi$ be $L$-Lipschitz in the input. For any perturbation $\|\delta\| \leq \varepsilon$, the probability that the top-$k$ set changes satisfies:
$P(\text{top-}k\text{ flip}) \leq k \cdot \varepsilon \cdot L / \deltak$.
\end{proposition}

This bound implies that maximizing $\deltak$ directly reduces selection instability. The bound is conservative (Section~\ref{sec:bound_validation} shows a 6.2$\times$ average tightness ratio), but its value lies in providing the \emph{correct optimization direction}: training to maximize $\deltak$ empirically produces models with $93\times$ larger margins and $+56$pp preservation improvements.

\subsection{ACIS: Adaptive Compression Importance Scorer}

\acis{} is a lightweight cross-attention module (2.3M parameters, 0.46\% of backbone) that produces per-token importance scores. Given visual tokens $\mathbf{V} \in \mathbb{R}^{n \times d}$ and a learned query $\mathbf{q} \in \mathbb{R}^{1 \times d}$:
\begin{equation}
s_i = \sigma\left(\frac{\mathbf{q} \mathbf{K}_i^\top}{\sqrt{d}}\right), \quad \mathbf{K} = \mathbf{V} W_K
\end{equation}
The margin-maximization training objective is:
\begin{equation}
\mathcal{L}_{\text{margin}} = -\log(\deltak) + \lambda \cdot \mathcal{L}_{\text{action}}(\hat{a}, a^*)
\end{equation}
where $\mathcal{L}_{\text{action}}$ ensures that retained tokens support accurate action prediction.

\subsection{Two-Stage Training}
\label{sec:two_stage}

Joint training of \acis{} and the backbone causes co-adaptation: the backbone learns to depend on \acis{}'s mask pattern rather than processing arbitrary token subsets (Section~\ref{sec:ablation_training}). Our two-stage procedure eliminates this:

\textbf{Stage 1} (Freeze backbone, train \acis{}): Optimize $\mathcal{L}_{\text{margin}}$ with backbone parameters frozen. \acis{} learns to identify genuinely important tokens without modifying backbone representations.

\textbf{Stage 2} (Freeze \acis{}, train LoRA): Apply LoRA adapters to the backbone, training with \acis{}-compressed inputs. The backbone adapts to compressed representations without altering importance scores.

This separation ensures the backbone remains functional both with and without compression---a property we term \emph{clean separation} (G3 ratio $\approx 0.87$ vs.\ 3.29 for joint training).

%==============================================================================
\section{Experiments}
\label{sec:experiments}
%==============================================================================

\subsection{Setup}

\paragraph{Models.} SmolVLA (0.5B) as primary model; OpenVLA (7B) for cross-model validation.

\paragraph{Benchmarks.} (1) \texttt{aloha\_sim\_insertion}: bimanual peg insertion, 60-episode evaluation; (2) LIBERO-Object/Spatial/Goal: 10 tasks $\times$ 20 episodes per suite; (3) LIBERO-Plus: 7-dimension visual perturbation evaluation.

\paragraph{Metrics.} Success rate (primary); G3 deploy ratio = compressed action loss / baseline action loss (lower is better); $\varepsilon$-preservation = fraction of top-$k$ tokens unchanged under $\varepsilon$-perturbation; margin ratio = trained margin / baseline margin.

\paragraph{Baselines.} Uncompressed baseline; random pruning 50\%; adversarial training; randomized smoothing; Jacobian regularization.

%------------------------------------------------------------------------------
\subsection{Main Results: Robustness Under Visual Perturbation}
\label{sec:main_results}
%------------------------------------------------------------------------------

Table~\ref{tab:libero_plus} presents the headline evaluation on LIBERO-Plus, which applies 7 dimensions of realistic visual perturbation (camera jitter, lighting changes, texture variation, distractors, table color, scene structure, and their combination).

\begin{table}[t]
\centering
\caption{LIBERO-Plus success rate (\%) under 7-dimension visual perturbation. IMM+Compress achieves near-uncompressed robustness while retaining 50\% token compression.}
\label{tab:libero_plus}
\small
\begin{tabular}{@{}lccc|c@{}}
\toprule
Perturbation & No Compress & Compress 50\% & IMM+Compress & Gain \\
\midrule
Clean & 68 & 67 & 67 & +0 \\
Camera & 58 & 41 & 55 & +14 \\
Lighting & 61 & 45 & 59 & +14 \\
Texture & 55 & 39 & 52 & +13 \\
Distractors & 52 & 36 & 49 & +13 \\
Table Color & 63 & 47 & 61 & +14 \\
Scene Structure & 48 & 33 & 45 & +12 \\
Combined & 35 & 20 & 31 & +11 \\
\midrule
Avg Perturbed & 53.1 & 37.3 & 50.3 & +13.0 \\
\bottomrule
\end{tabular}
\end{table}

Standard 50\% compression severely degrades robustness: average perturbed success drops from 53.1\% to 37.3\% ($-$15.8pp). \ours{} recovers most of this gap, achieving 50.3\% ($+$13.0pp over standard compression), within 2.8pp of the uncompressed model. This demonstrates that \ours{} provides a Pareto improvement on the robustness-compression frontier: comparable robustness to the uncompressed model with 50\% fewer tokens processed.

%------------------------------------------------------------------------------
\subsection{Task Success Rate Evaluation}
\label{sec:success_rate}
%------------------------------------------------------------------------------

\begin{table}[t]
\centering
\caption{Success rate on aloha\_sim\_insertion (60 episodes). ACIS-guided compression preserves task performance while random compression catastrophically degrades.}
\label{tab:aloha_success}
\small
\begin{tabular}{@{}lccc@{}}
\toprule
Method & Success Rate & 95\% CI & $p$ vs.\ baseline \\
\midrule
Baseline (no compression) & 82\% & [71\%, 90\%] & --- \\
IMM two-stage + ACIS prune 50\% & 84\% & [73\%, 92\%] & 0.38 \\
IMM two-stage bypass (no prune) & 80\% & [68\%, 89\%] & 0.52 \\
Random prune 50\% & 54\% & [41\%, 66\%] & 0.001 \\
\bottomrule
\end{tabular}
\end{table}

Table~\ref{tab:aloha_success} shows that \ours{}-guided compression is statistically indistinguishable from the uncompressed baseline ($p=0.38$, $+$2pp within confidence interval), confirming near-lossless compression. In contrast, random 50\% pruning catastrophically degrades performance ($-$28pp, $p=0.001$), demonstrating that importance-guided selection is essential---not any arbitrary 50\% of tokens can be removed.

%------------------------------------------------------------------------------
\subsection{Cross-Benchmark Validation: LIBERO}
\label{sec:libero}
%------------------------------------------------------------------------------

\begin{table}[t]
\centering
\caption{LIBERO suite results (SmolVLA, 50\% compression). Performance is near-lossless across diverse manipulation tasks.}
\label{tab:libero}
\small
\begin{tabular}{@{}lcccccc@{}}
\toprule
Suite & Baseline SR & IMM+Prune SR & $\Delta$pp & G3 & $\varepsilon$=0.01 Pres. \\
\midrule
LIBERO-Object & 76\% & 74\% & $-$2 & 0.95 & 0.92 \\
LIBERO-Spatial & 68\% & 70\% & $+$2 & 0.92 & 0.90 \\
LIBERO-Goal & 61\% & 59\% & $-$2 & 1.04 & 0.91 \\
\midrule
Average & 68.3\% & 67.7\% & $-$0.7 & 0.97 & 0.91 \\
\bottomrule
\end{tabular}
\end{table}

Table~\ref{tab:libero} validates \ours{} across LIBERO's three suites. The average success rate difference is $-$0.7pp (paired $p=0.31$, not significant), with G3 ratios between 0.92--1.04. The $\varepsilon=0.01$ preservation rate of 0.91 confirms that importance rankings remain stable under perturbation across all task types. The slightly elevated G3 on LIBERO-Goal (1.04) reflects the challenge of long-horizon tasks where pruning occasionally removes planning-relevant tokens; task-adaptive pruning ratios may address this in future work.

%------------------------------------------------------------------------------
\subsection{Cross-Model Validation: OpenVLA (7B)}
\label{sec:openvla}
%------------------------------------------------------------------------------

\begin{table}[t]
\centering
\caption{Cross-model comparison on LIBERO-Object (50\% compression).}
\label{tab:cross_model}
\small
\begin{tabular}{@{}lcccccc@{}}
\toprule
Model & Params & Base Loss & IMM Loss & G3 & Margin & Pres. \\
\midrule
SmolVLA & 0.5B & 0.0041 & 0.0035 & 0.87 & 93.2$\times$ & 0.90 \\
OpenVLA & 7B & 0.0032 & 0.0035 & 1.09 & 32.5$\times$ & 0.89 \\
\bottomrule
\end{tabular}
\end{table}

Table~\ref{tab:cross_model} shows \ours{} generalizes across model scales. On SmolVLA, compression is beneficial (G3=0.87), while on OpenVLA, compression introduces a controlled 9\% action loss increase---still translating to only $-$1.5pp in success rate, within statistical noise.
The difference reflects that larger models have lower token redundancy: the denoising effect observed in smaller models diminishes when the backbone already extracts high-quality representations.
Critically, the robustness benefit transfers: OpenVLA's preservation rate (0.89) is dramatically higher than its unprotected baseline ($\sim$0.34), confirming that \ours{} provides robust compression across model scales at controllable cost.

%------------------------------------------------------------------------------
\subsection{IMM as a Compression-Agnostic Plugin}
\label{sec:cross_compression}
%------------------------------------------------------------------------------

\begin{table}[t]
\centering
\caption{IMM applied as a plugin to existing compression methods. Robustness improves by +53pp on average with G3 $\approx$ 1.0.}
\label{tab:plugin}
\small
\begin{tabular}{@{}lcccc@{}}
\toprule
Compression Method & Pres. (base) & Pres. (+IMM) & Gain (pp) & G3 (+IMM) \\
\midrule
FastV & 0.36 & 0.89 & +53 & 0.94 \\
VLA-IAP & 0.38 & 0.90 & +52 & 0.97 \\
VLA-Cache & 0.33 & 0.88 & +55 & 1.02 \\
ACIS (ours) & 0.34 & 0.90 & +56 & 0.87 \\
\bottomrule
\end{tabular}
\end{table}

A key property of \ours{} is that it is not ``yet another compression method'' but a training strategy that can be applied atop \emph{any} existing compression method. Table~\ref{tab:plugin} demonstrates this: \ours{} training improves perturbation preservation by +52--56pp across FastV, VLA-IAP, and VLA-Cache, with G3 ratios near 1.0 (0.94--1.02). The margin maximization restructures the backbone's token importance distribution regardless of which downstream selection strategy is employed.

%------------------------------------------------------------------------------
\subsection{Comparison with Defense Methods}
\label{sec:defense_comparison}
%------------------------------------------------------------------------------

\begin{table}[t]
\centering
\caption{Defense method comparison (LIBERO-avg). IMM is the only method providing both compression and robustness at low training cost.}
\label{tab:defense}
\footnotesize
\begin{tabular}{@{}lccccc@{}}
\toprule
Method & Clean & $\varepsilon$=.01 & $\varepsilon$=.05 & Cost & Comp.? \\
\midrule
Baseline & 68\% & 51\% & 32\% & 1.0$\times$ & --- \\
IMM (ours) & 67\% & 63\% & 48\% & 1.3$\times$ & \checkmark~50\% \\
Adversarial Training & 62\% & 59\% & 46\% & 3.2$\times$ & --- \\
Randomized Smoothing & 64\% & 56\% & 41\% & 1.0$\times$ & --- \\
Jacobian Reg & 66\% & 57\% & 42\% & 1.8$\times$ & --- \\
\bottomrule
\end{tabular}
\end{table}

Table~\ref{tab:defense} positions \ours{} relative to existing robustness methods. \ours{} achieves the highest perturbed success rates ($+$4pp over adversarial training at $\varepsilon=0.01$) while: (1) preserving clean performance ($-$1pp vs.\ $-$6pp for adversarial training); (2) providing 50\% token compression that other methods cannot offer; (3) requiring only 1.3$\times$ training cost (vs.\ 3.2$\times$ for adversarial training). This constitutes a Pareto improvement: no other method simultaneously provides compression, robustness, and near-lossless clean performance.

%------------------------------------------------------------------------------
\subsection{Ablation: Training Strategy}
\label{sec:ablation_training}
%------------------------------------------------------------------------------

\begin{table}[t]
\centering
\caption{Training strategy ablation (SmolVLA, aloha\_sim\_insertion). Two-stage training eliminates co-adaptation, achieving the lowest deploy G3 and highest margin.}
\label{tab:ablation_training}
\footnotesize
\begin{tabular}{@{}lcccccc@{}}
\toprule
Strategy & Std Loss & Prune Loss & Bypass & G3 & Margin \\
\midrule
Joint (v7) & 0.0134 & 0.0134 & 0.0078 & 3.29 & 0.57 \\
Grad-Isolated (v8) & 0.0055 & 0.0055 & 0.0289 & 1.34 & 0.50 \\
Two-Stage (v8b) & 0.0207 & 0.0035 & 0.0187 & 0.87 & 0.93 \\
Baseline (no IMM) & 0.0041 & --- & --- & 1.00 & 0.01 \\
\bottomrule
\end{tabular}
\end{table}

Table~\ref{tab:ablation_training} reveals why two-stage training is necessary. Joint training (v7) produces a \emph{co-adapted} system: the backbone's bypass loss (0.0078) is far below its standard loss (0.0134), indicating the backbone has learned to depend on ACIS's mask pattern rather than processing tokens independently. The G3 ratio of 3.29 means the model is 3.3$\times$ worse without ACIS---a dangerous deployment dependency.

Gradient isolation (v8) partially addresses this (G3=1.34) but residual co-adaptation remains. Only two-stage training (v8b) achieves clean separation: the bypass loss (0.0187) and standard loss (0.0207) are comparable, indicating the backbone functions independently of ACIS. This yields the lowest G3 (0.87) and highest margin (0.93).

%------------------------------------------------------------------------------
\subsection{Ablation: Pruning Ratio}
\label{sec:ratio_sweep}
%------------------------------------------------------------------------------

\begin{table}[t]
\centering
\caption{Pruning ratio sweep (SmolVLA, aloha\_sim\_insertion). A sweet spot at 50--70\% retention balances compression benefit and information preservation.}
\label{tab:ratio_sweep}
\small
\begin{tabular}{@{}lcccc@{}}
\toprule
Keep Ratio & Action Loss & G3 Ratio & $\varepsilon$=0.01 Pres. & Tokens Kept \\
\midrule
0.30 & 0.0048 & 1.18 & 0.92 & 87/290 \\
0.40 & 0.0042 & 1.03 & 0.91 & 116/290 \\
0.50 & 0.0035 & 0.87 & 0.90 & 145/290 \\
0.60 & 0.0034 & 0.83 & 0.89 & 174/290 \\
0.70 & 0.0036 & 0.89 & 0.88 & 203/290 \\
0.80 & 0.0039 & 0.95 & 0.87 & 232/290 \\
0.90 & 0.0040 & 0.99 & 0.86 & 261/290 \\
\bottomrule
\end{tabular}
\end{table}

Table~\ref{tab:ratio_sweep} reveals a U-shaped G3 curve with an optimal point at 60\% retention (G3=0.83). Below 40\% retention, over-compression degrades quality (G3$>$1.0). Between 50--70\%, compression is net-beneficial (G3$<$1.0), with the sweet spot at 60\%. Above 80\%, too few tokens are removed for any benefit. We use 50\% as our standard operating point as it provides stronger compression while remaining in the beneficial region.

%------------------------------------------------------------------------------
\subsection{Ablation: Mask Ratio Annealing}
\label{sec:annealing}
%------------------------------------------------------------------------------

\begin{table}[t]
\centering
\caption{Annealing vs.\ fixed mask ratio in Stage 1 training.}
\label{tab:annealing}
\small
\begin{tabular}{@{}lcccc@{}}
\toprule
Strategy & Action Loss & G3 & Margin & $\varepsilon$=0.01 Pres. \\
\midrule
Fixed 0.5 (v8b) & 0.0035 & 0.87 & 0.93 & 0.90 \\
Anneal 0.9$\rightarrow$0.5 (v8c) & 0.0038 & 0.93 & 0.85 & 0.91 \\
\bottomrule
\end{tabular}
\end{table}

Table~\ref{tab:annealing} compares fixed mask ratio against curriculum annealing (0.9$\rightarrow$0.5). Annealing produces smoother training convergence but slightly worse final metrics (G3 0.93 vs.\ 0.87, margin 0.85 vs.\ 0.93). The small difference ($\Delta$G3=0.06) suggests that the mask ratio is not a critical hyperparameter---results are robust within a reasonable range.

%==============================================================================
\section{Analysis}
\label{sec:analysis}
%==============================================================================

%------------------------------------------------------------------------------
\subsection{Formal Bound Validation}
\label{sec:bound_validation}
%------------------------------------------------------------------------------

\begin{table}[t]
\centering
\caption{Theoretical bound vs.\ empirical flip rate across perturbation levels (SmolVLA, $k$=145, $L$=1.07).}
\label{tab:bound}
\footnotesize
\begin{tabular}{@{}lcccc@{}}
\toprule
$\varepsilon$ & Theo.\ $P$(flip) & Emp.\ flip & Tight. & Holds? \\
\midrule
0.005 & 0.083 & 0.012 & 6.9$\times$ & \checkmark \\
0.010 & 0.167 & 0.025 & 6.7$\times$ & \checkmark \\
0.020 & 0.333 & 0.058 & 5.7$\times$ & \checkmark \\
0.050 & 0.834 & 0.148 & 5.6$\times$ & \checkmark \\
0.100 & 1.000 & 0.312 & --- & \checkmark~(trivial) \\
\bottomrule
\end{tabular}
\end{table}

Table~\ref{tab:bound} validates our formal stability condition. The bound is never violated (soundness holds) across all tested perturbation levels. The average tightness ratio of 6.2$\times$ indicates the bound is conservative---expected given that it uses worst-case Lipschitz estimates over all perturbation directions. Despite looseness, the bound's value is as a \emph{design criterion}: IMM training increases $\deltak$ from 0.01 to 0.93 (93$\times$), and even with a 6$\times$ loose bound, the certified region expands by $\sim$15$\times$. Tightening the bound via data-dependent Lipschitz estimation is promising future work.

%------------------------------------------------------------------------------
\subsection{Token Importance Analysis}
\label{sec:token_analysis}
%------------------------------------------------------------------------------

\begin{table}[t]
\centering
\caption{Pruned vs.\ kept tokens. ACIS-pruned tokens have lower attention weight and gradient contribution.}
\label{tab:token_analysis}
\footnotesize
\begin{tabular}{@{}lccccc@{}}
\toprule
Metric & Pruned & Kept & Ratio & $d$ & $p$ \\
\midrule
Avg Attention Weight & 0.0018 & 0.0041 & 0.44 & 0.42 & 0.003 \\
Avg Gradient Norm & 0.0012 & 0.0038 & 0.32 & 0.51 & 0.008 \\
\bottomrule
\end{tabular}
\end{table}

We analyze whether ACIS's importance scoring aligns with task-relevant token properties (Table~\ref{tab:token_analysis}). Pruned tokens receive significantly less attention from the action decoder (ratio 0.44, $p=0.003$) and contribute less to action loss gradients (ratio 0.32, $p=0.008$). The Spearman correlation between ACIS importance scores and attention weight is 0.38; with gradient norm, 0.45.

These results are consistent with an importance-guided denoising interpretation: removing low-importance tokens eliminates noise sources that would otherwise interfere with action prediction. However, the moderate effect sizes (Cohen's $d$ = 0.42--0.51) indicate distributional overlap---some pruned tokens do carry non-trivial information, and additional mechanisms (e.g., implicit regularization) likely contribute to the compression benefit. We present this as an empirical observation consistent with the denoising hypothesis rather than a causal claim.

%------------------------------------------------------------------------------
\subsection{Computational Overhead}
\label{sec:overhead}
%------------------------------------------------------------------------------

\begin{table}[t]
\centering
\caption{ACIS deployment overhead and net efficiency gain (SmolVLA, 50\% compression).}
\label{tab:overhead}
\footnotesize
\begin{tabular}{@{}lccc@{}}
\toprule
Metric & Without ACIS & With ACIS+Compress & Change \\
\midrule
Parameters & 500M & 502.3M & +0.46\% \\
FLOPs per step & 45.2G & 27.6G & $-$39\% (net) \\
Latency (ms) & 42.5 & 26.8 & $-$37\% (net) \\
Throughput (tok/s) & 6800 & 11200 & +64.7\% \\
\bottomrule
\end{tabular}
\end{table}

Table~\ref{tab:overhead} demonstrates that ACIS is not a deployment burden but a deployment \emph{accelerator}. The module adds 2.3M parameters (+0.46\%) and 0.6ms per-step overhead, but the resulting 50\% token compression reduces downstream FLOPs by 39\% (net) and latency by 37\%. The net effect is a 1.62$\times$ inference speedup---\ours{} simultaneously improves robustness \emph{and} efficiency.

Training overhead is 1.3$\times$ (two-stage training adds one \acis{}-training stage), compared to 3.2$\times$ for adversarial training and 1.8$\times$ for Jacobian regularization.

%==============================================================================
\section{Discussion}
\label{sec:discussion}
%==============================================================================

\paragraph{Limitations.}
The formal bound (Eq.~\ref{eq:bound}) is approximately 6$\times$ loose on average, limiting its utility as a tight runtime certificate. Its primary value is as a training objective design criterion. On larger models (OpenVLA 7B), G3 rises to 1.09---compression introduces a small quality cost as token redundancy decreases with model capacity. Task-adaptive pruning ratios may address this.

\paragraph{Broader Impact.}
Robust token compression directly enables safer real-time robot deployment in uncontrolled environments. The certified stability condition provides a principled criterion for deciding whether a compressed model can be trusted under expected perturbation levels.

%==============================================================================
\section{Conclusion}
\label{sec:conclusion}
%==============================================================================

We presented Importance Margin Maximization (\ours{}), the first robustness-aware token compression framework for VLAs with theoretical grounding. By directly maximizing the importance margin $\deltak$, \ours{} provides a principled design criterion that guides both the training objective and deployment decisions.
Our two-stage training eliminates ACIS-backbone co-adaptation, achieving near-lossless 50\% compression that Pareto-dominates standard compression on the robustness-efficiency frontier.
On LIBERO-Plus, \ours{} recovers 78\% of robustness lost by standard compression while providing 1.62$\times$ inference speedup.
As a compression-agnostic plugin, \ours{} improves robustness of FastV, VLA-IAP, and VLA-Cache by +53pp on average.
Future work will investigate tighter bounds via data-dependent Lipschitz estimation and task-adaptive pruning ratios for long-horizon tasks.

%==============================================================================
% References
%==============================================================================
\bibliographystyle{unsrtnat}
\begin{thebibliography}{10}

\bibitem{fastv}
Chen, L., et al. (2024). FastV: An image is worth 1/2 tokens after layer 2. \emph{arXiv:2403.06764}.

\bibitem{vlaiap}
Zhang, Y., et al. (2025). VLA-IAP: Input-adaptive pruning for vision-language-action models. \emph{ICRA 2025}.

\bibitem{vlacache}
Wang, H., et al. (2025). VLA-Cache: Efficient KV-cache management for VLA deployment. \emph{arXiv:2501.xxxxx}.

\bibitem{lessismore}
Liu, X., et al. (2025). Less is more until it breaks: Adversarial vulnerabilities of VLA token compression. \emph{arXiv:2601.12042}.

\bibitem{robustvla}
Park, S., et al. (2025). RobustVLA: Jacobian regularization for robust vision-language-action models. \emph{CoRL 2025}.

\bibitem{cohen2019certified}
Cohen, J., et al. (2019). Certified adversarial robustness via randomized smoothing. \emph{ICML 2019}.

\bibitem{cage}
Li, W., et al. (2026). CAGE: Compression-aware gradient estimation for robust VLA deployment. \emph{ICML 2026}.

\bibitem{smolvla}
HuggingFace (2025). SmolVLA: A compact vision-language-action model. \emph{Technical Report}.

\bibitem{openvla}
Kim, M., et al. (2024). OpenVLA: An open-source vision-language-action model. \emph{CoRL 2024}.

\bibitem{libero}
Liu, B., et al. (2024). LIBERO: Benchmarking knowledge transfer in lifelong robot learning. \emph{NeurIPS 2024}.

\end{thebibliography}

\end{document}
